\chapter{Construcciones Geométricas Básicas en CAD}

\section{Fundamentos de Construcciones Geométricas}

Las construcciones geométricas constituyen el lenguaje fundamental del
dibujo técnico. En CAD, estas operaciones se realizan con una
precisión inigualable, permitiendo transformar primitivas simples en
diseños mecánicos o arquitectónicos complejos.

\subsection{Primitivas Geométricas}

Los programas CAD utilizan entidades básicas:
\begin{itemize}
	\item \textbf{Líneas y Polilíneas:} Son los elementos base. Las
	      polilíneas permiten crear un objeto compuesto de múltiples
	      segmentos rectos y curvos tratados como una sola entidad.
	\item \textbf{Círculos y Arcos:} Definidos habitualmente por su
	      centro y radio, o por puntos de paso (inicio, fin, radio).
	\item \textbf{Elipses:} Definidas por sus ejes mayor y menor.
\end{itemize}

\subsection{Herramientas de Precisión}

La verdadera potencia del CAD reside en sus herramientas de ayuda:
\begin{itemize}
	\item \textbf{OSNAP (Objetos de ajuste o Referencias a objetos):}
	      Es una herramienta esencial que permite al cursor "pegarse"
	      a puntos geométricos clave de las entidades existentes, como
	      el centro de un círculo, el punto medio de una línea, una
	      intersección o un punto final.
	\item \textbf{Modo Orto:} Fuerza el movimiento del cursor para
	      que las líneas solo se dibujen de forma horizontal o
	      vertical (ángulos de 0, 90, 180, 270 grados).
	\item \textbf{Rastreo Polar:} Permite dibujar líneas según
	      ángulos prefijados, facilitando construcciones inclinadas
	      precisas.
\end{itemize}

\subsection{Operaciones de Modificación}

Además de dibujar, el CAD ofrece comandos para modificar geometrías:
\begin{itemize}
	\item \textbf{Empalme (Fillet):} Crea un arco tangente entre dos
	      entidades (líneas, arcos).
	\item \textbf{Chaflán (Chamfer):} Crea una línea inclinada (bisel)
	      entre dos entidades.
	\item \textbf{Recortar/Extender (Trim/Extend):} Modifica las
	      longitudes de las entidades basadas en otras entidades
	      cercanas.
\end{itemize}

La tangencia es una propiedad crítica en el diseño mecánico, asegurando
transiciones suaves entre curvas y rectas, lo cual es fundamental para
la correcta funcionalidad de piezas móviles y ensamblajes.

\section{Cuestionario}
\begin{enumerate}
	\item ¿Qué es el modo Orto y en qué situaciones es útil?
	\item ¿Cómo ayuda el OSNAP a la precisión del dibujo?
	\item Define qué es un círculo, radio y diámetro.
	\item ¿Cómo se traza un arco tangente a dos líneas?
	\item ¿Qué son los puntos de control en una curva?
	\item Explica la importancia de la tangencia en el diseño mecánico.
	\item ¿Cómo se calcula el área de una figura compleja mediante CAD?
	\item ¿Qué es el perímetro?
\end{enumerate}

\section{Parte Práctica}
\begin{enumerate}
	\item Dibujar un triángulo equilátero de 50 mm de lado.
	\item Dibujar un polígono regular de 6 lados inscrito en un círculo de 40 mm de radio.
	\item Trazar dos círculos y unirlos mediante una línea tangente.
	\item Dividir una línea de 100 mm en 5 partes iguales.
	\item Dibujar un arco de 60 grados con radio 30 mm.
	\item Construir una elipse con radios 50 y 30 mm.
	\item Dibujar un cuadrado de 40 mm y redondear sus vértices con radio 5 mm.
	\item Trazar dos líneas perpendiculares y unirlas con un arco de radio 10 mm.
	\item Dibujar un círculo tangente a tres líneas dadas.
	\item Crear una figura compleja compuesta por arcos y líneas.
\end{enumerate}
